% =====================================================================
%  One Stack, Two Ledgers
%  CollusionGraph -- research paper, written from docs/paper_blueprint.md
%
%  COMPILING ON OVERLEAF
%  ---------------------
%  Upload this .tex together with the seven PNG figures that live in the
%  repository root:
%      fig1_architecture.png        fig5_label_noise.png
%      fig2_two_ledgers.png         fig6_label_efficiency.png
%      fig3_timestep_crater.png     fig7_dual_actor.png
%      fig4_transfer_lift.png
%  Everything else is self-contained -- the tables are inlined, so there
%  is nothing to \input and no external .bib to resolve.
%
%  HOUSE RULE: every number in this paper is traced to an artifact under
%  eval_outputs/ via docs/reproducibility.md. Regenerate the tables with
%  `uv run poe paper-tables` before editing any figure quoted here.
% =====================================================================
\documentclass[conference]{IEEEtran}
\IEEEoverridecommandlockouts

\usepackage{cite}
\usepackage{amsmath,amssymb}
\usepackage{graphicx}
\usepackage{booktabs}
\usepackage{array}
\usepackage{xcolor}
\usepackage{balance}
\usepackage[hidelinks]{hyperref}
\usepackage{url}

\newcommand{\RQ}[1]{\textbf{RQ#1}}

\begin{document}

\title{One Stack, Two Ledgers: Budget-Evaluated, Explainable\\
Collusion Screening Across Illicit-Finance and\\
Procurement Networks}

\author{%
\IEEEauthorblockN{Kartik Joshi, Prem Kukreja, Gagandeep Singh, Samuel Alex}
\IEEEauthorblockA{\textit{[INSTITUTION -- replace before submission]}\\
\{kartik, prem, gagandeep, samuel\}@example.edu}
\and
\IEEEauthorblockN{Rachit Garg}
\IEEEauthorblockA{\textit{[INSTITUTION -- replace before submission]}\\
rachit.garg@example.edu}
}

\maketitle

\begin{abstract}
Money laundering and bid rigging are studied by separate literatures, but
both are \emph{structural} crimes: no single transaction or tender is
incriminating, and the evidence lives in the shape of the network. We build
one graph-learning stack --- a single intermediate representation, one
detector family, one explanation layer and one evaluation protocol --- and
exercise it on six public datasets across both domains. Under a strict
inductive, budget-first protocol, gradient-boosted trees on graph features
still beat every graph neural network we trained on Elliptic++ (XGB-Graph
$0.8104$ AUC-PR against GATv2-focal $0.4729 \pm 0.0525$ over five seeds;
paired bootstrap $\Delta = -0.261$, 95\% CI $[-0.285, -0.235]$,
$p \approx 0.001$): the GADBench finding replicates, and multi-seed
evaluation sharpens rather than softens it. Transfer across procurement
markets is real but \emph{market-dependent} --- leave-one-market-out on
Garc\'ia is uniformly positive (macro lift $1.57$) while leave-one-country-out
on the Mendeley EU cartel data averages a lift of $1.17$ yet \emph{fails on
its largest market} ($0.90$). Cross-domain transfer is asymmetric and
label-rich-only: a frozen procurement encoder lifts financial screening
$2.31\times$ over prevalence, the reverse direction scores below prevalence,
and neither pays at fewer than roughly a thousand target labels. We report a
sharp honest negative: corrupting 20\% of training labels \emph{raises} test
AUC-PR from $0.483$ to $0.598$ while validation collapses from $0.942$ to
$0.569$, making temporal-shift validation-blindness explicit. Alongside the
measurements we ship the product: PGExplainer subgraphs with motif matching
and FATF/OECD red-flag citations, a fixed-budget investigator console, and a
guarded LLM copilot that passes a 24-case golden gate with perfect grounding
and zero guilt-attributing releases. Code, configs, seeds, five datasheets, a
model card and a number-to-config map are released. The system is a screening
and triage instrument; it makes no determination of guilt.
\end{abstract}

\begin{IEEEkeywords}
graph neural networks, anomaly detection, anti-money laundering, bid rigging,
class imbalance, explainability, transfer learning, evaluation protocol
\end{IEEEkeywords}

% =====================================================================
\section{Introduction}
% =====================================================================

On 2 August 2026 the substantive obligations of the EU AI Act become
enforceable for high-risk systems, a class that explicitly includes fraud
detection in financial services. Explainability, human oversight, logging and
auditability stop being product niceties and become legal properties of the
system. Any screening model deployed after that date must be able to say
\emph{why} it raised a flag, in terms a human reviewer can check and a
regulator can audit.

That deadline lands on two enforcement worlds that have never shared a
toolchain. Anti-money-laundering (AML) teams inside banks run
transaction-monitoring rules that fire on individual records; industry
benchmarks put the false-positive share of those alerts at 90--95\%.
Competition authorities and procurement auditors run statistical screens over
tenders one at a time, by hand, on a small sample, and typically discover
cartels only afterwards, through complaints or leniency applications.

The two worlds are closer than their literatures suggest. A single wire
transfer is unremarkable; twelve accounts passing the same funds around a
circle within 72 hours is laundering. A single bid is procedurally clean --- the
tender was advertised, bids were received, a winner was declared; five firms
that quietly rotate the winner among themselves, submit implausibly clustered
losing prices, and share directors and registered addresses are a cartel. In
both cases the individual record is innocent and \emph{the network is guilty}.
In both cases the participants deliberately keep every record below the
threshold the rules check, so the illegality is only visible one level up, in
the topology. And in both cases confirmed positives are vanishingly rare
relative to clean activity.

Fig.~\ref{fig:ledgers} makes the correspondence concrete. Five collusion
mechanisms recur across the domains, and each reduces to the same graph
signature: laundering cycles and bid rotation are both directed cycles;
smurfing and cover bidding are both fan-in stars; layering and market
allocation are both fan-out stars; shell-company chains and shared
directorships are both dense co-attribute subgraphs; pass-through clusters and
clustered bid prices are both tight communities with anomalous edge
attributes. If that correspondence is real, one stack should serve both
ledgers --- and whether a learned notion of ``collusion structure'' carries
from one to the other becomes an empirical question rather than a slogan.

\begin{figure*}[t]
\centering
\includegraphics[width=\textwidth]{fig2_two_ledgers.png}
\caption{Two ledgers, one structure. Five collusion mechanisms, the graph
signature each leaves behind, and the domain-specific name it goes by on
either side. The correspondence is the premise of the shared stack: only the
adapters and the red-flag vocabularies differ below it.}
\label{fig:ledgers}
\end{figure*}

\subsection{Research questions}

\begin{itemize}
\item[\RQ{1}] Can graph neural networks combined with unsupervised graph
anomaly detection identify illicit-finance and bid-rigging structures with
materially higher precision \emph{at fixed alert budgets} than rule-based and
tabular machine-learning baselines?
\item[\RQ{2}] Do collusion fingerprints learned from confirmed cases,
augmented with synthetically injected patterns, transfer to unseen markets and
unseen typologies at a controlled flag rate?
\item[\RQ{3}] Can subgraph-level explanation methods produce alert rationales
that align with the red-flag frameworks practitioners already use?
\item[\RQ{4}] To what extent can one shared stack serve both domains, and does
a representation of collusion structure learned in one provide usable signal
in the other?
\end{itemize}

\subsection{Contributions}

\begin{enumerate}
\item \textbf{A unified intermediate representation and six adapters} spanning
both domains, with degradation paths (missing bidder identities, case-control
sampling, post-window label spikes) recorded rather than repaired.
\item \textbf{A budget-first, leakage-safe evaluation protocol} with strict
inductive splits, frozen train-time normalisation, tie-aware precision@$k$,
multi-seed reporting and paired stratified bootstrap significance. We treat
the protocol itself as a contribution: recent critiques of graph-anomaly
evaluation are implemented here, not merely cited.
\item \textbf{Transfer matrices at both granularities and in both
directions}, including label-efficiency curves that locate where cross-domain
transfer stops paying.
\item \textbf{An explanation layer with an instrument attached} --- PGExplainer
subgraphs, a nine-family motif matcher, FATF/OECD citations, a practitioner
study protocol, and a guarded copilot with a released-output golden gate.
\item \textbf{A catalogue of honest negatives}, including four independent
measurements that validation selection is blind to the Elliptic distribution
shift, and a joint cross-crime experiment specified but deliberately not yet
run (Section~\ref{sec:joint}).
\end{enumerate}

% =====================================================================
\section{Related Work}
% =====================================================================

\textbf{Graph anomaly detection and benchmarks.} Ma et al.'s survey gives the
unifying node/edge/subgraph taxonomy this work sits in. GADBench is the
uncomfortable reference point: across a wide sweep of datasets, well-tuned
tree ensembles over graph-derived features match or beat specialised graph
neural networks. Our \RQ{1} result replicates that finding on Elliptic++ under
a protocol strict enough that we consider the replication itself informative.

\textbf{AML with graph neural networks.} Motie and Raahemi survey the space;
LineMVGNN adds line-graph views for directed money flow; Lawal et al. apply
GNNExplainer to Elliptic and document that performance stays unstable under
class imbalance even with class weighting and focal loss --- the motivating
evidence that headline accuracy is deceptive at 2\% prevalence. The
Multi-GNN line reports that bidirectional message passing is decisive on
AMLworld; we replicate that in-house ($-0.194$ AUC-PR without reverse edges,
our strongest single component ablation). Reported near-perfect Elliptic
results should be read against their own ablations; at least one recent
wavelet-transformer paper concedes its distinguishing features contribute
about $0.01$.

\textbf{Cartel detection.} The screens tradition (coefficient of variation,
price spread, kurtosis) is well established, and Garc\'ia Rodr\'iguez et al.\
provide the reused multi-country benchmark. Recent cartel GNN work --- graph
attention networks over tender--firm graphs across 13 markets, and relational
GCNs with an explicit zero-shot phase --- reaches strong within-market numbers
but documents consistent cross-country degradation. Our leave-one-market-out
and leave-one-country-out matrices are the direct comparison point, and they
show the degradation is not uniform: it is market-dependent.

\textbf{Evaluation-protocol critiques.} A recent thread argues that
graph-anomaly results are inflated by leakage, transductive splits and
granularity mismatch between the scored unit and the operational unit. Our
transaction-versus-actor disagreement (Section~\ref{sec:granularity}) is a
direct measurement of the last of these.

\textbf{Transfer and graph foundation models.} The graph-foundation-model
framing asks whether one representation serves many graphs. \RQ{4} is a
falsifiable, negative-tolerant case study of exactly that claim across two
enforcement domains.

\textbf{LLM-augmented explanation.} Recent work pairs language models with
graph explanations to produce analyst-facing narratives. Our copilot is
deliberately bounded --- read-only SQL over published artifacts, a numeric
sanity gate, and a guilt-language guard --- so that it supports participation
without creating dependence.

% =====================================================================
\section{Data and Unified Representation}
% =====================================================================

\subsection{Datasets}

Table~\ref{tab:data} lists the six public datasets. Two are labelled anchors
that carry the headline results; the rest supply scale, controllable ground
truth, and an unlabelled population for injection studies. Every dataset is
downloaded from its canonical source and verified against a committed
checksum manifest that also records its licence. No raw data is redistributed
in the release.

\begin{table*}[t]
\caption{The six public datasets. Roles: \textsc{anchor} carries headline
results, \textsc{scale} stresses the pipeline, \textsc{ground truth} supplies
known injected patterns, \textsc{population} supplies an unlabelled portfolio.}
\label{tab:data}
\centering
\footnotesize
\setlength{\tabcolsep}{4pt}
\begin{tabular}{llllll}
\toprule
Dataset & Domain & Scale & Positives & Role & Licence \\
\midrule
Elliptic++ (tx graph) & Financial & 203,769 nodes, 234,355 edges, 49 steps & $\approx$2\% illicit & \textsc{anchor} & research use \\
Elliptic++ (actor graph) & Financial & 292,102 test nodes, 966,352 edges & 5.3\% prevalence & granularity & research use \\
IBM AMLworld HI-Small & Financial & 5,078,345 tx, 515,080 accounts & 0.102\% laundering & \textsc{scale} & CDLA-Sharing-1.0 \\
Mendeley EU cartel & Procurement & 15,616 contracts, 7 countries, 2004--21 & 41.9\% (case-control) & \textsc{anchor} & CC BY-NC 3.0 \\
Garc\'ia Rodr\'iguez et al. & Procurement & 64,348 bids, 9,781 tenders, 6 markets & 11.3--81.8\% by market & transfer & CC BY-NC-ND 4.0 \\
OCDS Georgia (OpenTender) & Procurement & 163,327-node window & unlabelled & \textsc{population} & CC BY-NC-SA 4.0 \\
\bottomrule
\end{tabular}
\end{table*}

\subsection{The intermediate representation}

Every adapter emits the same typed structure --- nodes, directed attributed
edges, labels with provenance and confidence, and a history pack for as-of
feature computation --- so that everything downstream is domain-agnostic. The
contract forbids cleaning heroics: an adapter may not impute a missing
identity, invent an edge, or silently drop a degradation. Three degradations
are therefore explicit in the representation rather than hidden:

\begin{itemize}
\item \textbf{Garc\'ia bidder identities exist on four of six markets.} The
company-identifier column is present only in the per-market files for Japan,
Italy, Brazil and the USA; the combined file drops it and the two Swiss
markets never carried it. Co-bidding \emph{identity} graphs are therefore
constructible on four markets; the remainder carry the anonymous bid-price
tier only.
\item \textbf{Mendeley is a case-control research sample}, not a population
file: 41.9\% of its contracts are cartel-linked. Raw AUC-PR is consequently
not comparable across its folds, whose prevalences range from $0.29$ to
$0.89$. \emph{Lift over fold prevalence, never raw AUC-PR, is the cross-fold
comparator} throughout this paper.
\item \textbf{AMLworld has a post-window label spike.} The 1,108 transactions
dated after 10 September are 59.1\% laundering, roughly 580 times the base
rate. The temporal splitter fences that tail explicitly; a naive split would
manufacture a spectacular and entirely artificial result.
\end{itemize}

\subsection{As-of label and feature policies}

Any feature computed from a node's history is computed \emph{as of} the split
boundary, never over the full record. This is a leakage contribution rather
than an implementation detail: history-derived features are exactly where
temporal leakage hides, because the leak is invisible in aggregate statistics
and shows up only as a suspiciously strong validation score. The policies are
enforced by a dedicated leakage test suite that runs as its own continuous
integration step on every push.

% =====================================================================
\section{Method}
% =====================================================================

Fig.~\ref{fig:arch} shows the whole system. The design rule is a hard trust
boundary: the offline pipeline writes plain files and exits; everything
user-facing is read-only and cannot retrain a model or alter a score.

\begin{figure*}[t]
\centering
\includegraphics[width=0.94\textwidth]{fig1_architecture.png}
\caption{The shared stack. Stages 1--7 run offline and terminate, writing
Parquet and JSON artifacts; stages 8--9 serve those artifacts read-only. The
same nine stages run for both domains --- only the ingestion adapters and the
red-flag vocabularies differ.}
\label{fig:arch}
\end{figure*}

\subsection{Detection arms}

Four arms score every node and are fused afterwards.

\textbf{Baselines.} B1 is a rule engine in the spirit of deployed
transaction-monitoring thresholds. B2 is XGBoost on node features. B3 is
XGBoost on node features \emph{plus} graph-derived features (degrees, k-core,
motif counts, neighbourhood aggregates) --- the GADBench discipline, and the
yardstick that matters. B4, on the procurement side only, is the classical
screens tier (coefficient of variation, price spread, kurtosis, relative
distance).

\textbf{Supervised graph neural networks.} GraphSAGE, GATv2 and R-GCN, all
direction-aware: the graph is built with explicit reverse edges so that
message passing can distinguish sending from receiving. Focal loss is the
default given roughly 2\% positives; weighted cross-entropy is the ablation
arm.

\textbf{Unsupervised arm.} DOMINANT and a graph autoencoder score structural
reconstruction error, plus a deterministic structural floor. These require no
labels and are the only arms available on the unlabelled OCDS population.

\textbf{Fusion.} Each member is isotonically calibrated on validation before
fusion, so that the combined score is a comparable probability rather than an
arbitrary rank blend. The contrast between calibrated and rank fusion is one
of the paper's largest measured effects.

\subsection{From node scores to an alert queue}

Investigator capacity, not model output, is the binding resource, so the
operational unit is a \emph{group}, not a node. Leiden community detection
over the scored graph produces candidate groups; groups are ranked by
aggregate member risk, deduplicated by Jaccard non-maximum suppression,
capped in size, and truncated to a fixed review budget. A group counts as a
hit under an explicit hit rule (at least one confirmed member, tie-aware).
Every parameter of this pipeline is swept in Section~\ref{sec:robustness}.

\subsection{Explanation layer}

For each alert the system returns the minimal responsible subgraph
(PGExplainer, adopted over GNNExplainer after the fidelity ablation in
Section~\ref{sec:rq3}), the matched motif from a nine-family matcher, the
supporting attributes, and the specific FATF or OECD red-flag indicator the
shape corresponds to. An alert without an explanation is treated as a
non-deliverable.

\subsection{Guarded copilot}

A bounded language-model agent answers analyst questions over the published
artifacts. It is constrained by construction: a SELECT-only allowlist over
read-only views, a numeric sanity gate that rejects any figure not present in
the underlying rows, a guilt-language guard, and an immutable caveat on every
response. Its release gate is a golden set of 24 cases graded for grounding
and for released-output violations.

\subsection{Evaluation protocol}
\label{sec:protocol}

This subsection is the differentiator, and every clause of it changed a
headline number at some point during development.

\begin{itemize}
\item \textbf{Strict inductive splits.} Test nodes and their edges are unseen
at training time. Temporal splits fence the boundary; entity splits (LOCO,
LOMO) isolate whole countries or markets.
\item \textbf{Frozen train-time normalisation.} Feature statistics are
computed on train and frozen. Recomputing them on test is a form of
inadvertent test-time adaptation; removing it dropped our own headline GATv2
number from $0.693$ to $0.532$ and widened the honest gap to the tree
baseline. We report the post-fix numbers throughout.
\item \textbf{Budget-first, tie-aware metrics.} Precision@$k$ at the budgets a
real team would work, with explicit tie handling, alongside AUC-PR quoted
against the prevalence baseline of its own fold.
\item \textbf{Multi-seed with dispersion.} Five seeds, reported as
mean~$\pm$~standard deviation. Seed variance ($\pm 0.05$) dominates
cross-machine variance ($\pm 0.02$), so a single-seed comparison across
machines is uninterpretable. \emph{The published seed-0 GATv2 number,
$0.5492$, is the best of five seeds; the honest figure is
$0.4729 \pm 0.0525$, and this paper quotes the mean.}
\item \textbf{Paired stratified bootstrap} over identical confirmed test
nodes for every headline comparison, 2,000 resamples.
\end{itemize}

% =====================================================================
\section{Results}
% =====================================================================

\subsection{RQ1: trees still win, and multi-seed sharpens it}

Table~\ref{tab:t1} gives the Elliptic++ scoreboard. Gradient-boosted trees on
graph features reach $0.8104$ AUC-PR against a prevalence baseline of
$0.065$. The best graph neural network we trained reaches
$0.4729 \pm 0.0525$. The gap is not a tuning artifact of one seed:
Table~\ref{tab:t3} puts the paired bootstrap delta at $-0.261$ with a 95\%
confidence interval of $[-0.285, -0.235]$ and $p \approx 0.001$.

\begin{table}[t]
\caption{Elliptic++ node-level results (test steps 35--49, prevalence
$0.065$). Rule and tree baselines are deterministic; GNN and ensemble rows are
5-seed mean $\pm$ standard deviation.}
\label{tab:t1}
\centering
\footnotesize
\setlength{\tabcolsep}{4pt}
\begin{tabular}{lll}
\toprule
Model & AUC-PR & P@100 \\
\midrule
B1 rules & 0.0576 (det.) & 0.00 \\
B2 XGBoost & 0.8076 (det.) & 1.00 \\
\textbf{B3 XGBoost + graph} & \textbf{0.8104} (det.) & \textbf{1.00} \\
GATv2-focal & $0.4729 \pm 0.0525$ & $0.8120 \pm 0.2383$ \\
GATv2-weighted CE & $0.4388 \pm 0.0505$ & $0.7240 \pm 0.1466$ \\
DOMINANT & $0.0410 \pm 0.0003$ & --- \\
Ensemble, calibrated & $0.4434 \pm 0.0501$ & --- \\
Ensemble, rank & $0.0511 \pm 0.0019$ & --- \\
Structural floor & $0.0547 \pm 0.0000$ & --- \\
GAE & $0.0386 \pm 0.0001$ & --- \\
\bottomrule
\end{tabular}
\end{table}

\begin{table}[t]
\caption{Paired-bootstrap comparisons over identical confirmed test nodes
(2,000 resamples, stratified; seed-0 score files, so these deltas measure
sampling uncertainty at fixed seed and are not differences of multi-seed
means).}
\label{tab:t3}
\centering
\footnotesize
\setlength{\tabcolsep}{4pt}
\begin{tabular}{lrrr}
\toprule
Comparison & $\Delta$ AUC-PR & 95\% CI & $p$ \\
\midrule
GATv2-focal vs.\ B3 & $-0.261$ & $[-0.285, -0.235]$ & 0.001 \\
Calibrated vs.\ rank & $+0.471$ & $[0.440, 0.499]$ & 0.001 \\
\bottomrule
\end{tabular}
\end{table}

Two things are worth separating. First, the tree baseline is not beating the
graph model because the graph is uninformative --- B3's margin over B2
($0.8104$ vs.\ $0.8076$) is small, so most of its power is in the node
features, and the deep models are simply failing to convert structure into
generalisation under this shift. Second, the largest effect in
Table~\ref{tab:t3} is not model choice at all: calibrating ensemble members
before fusing them is worth $+0.471$ AUC-PR over rank fusion, an order of
magnitude more than any architectural decision we tested.

Fig.~\ref{fig:crater} explains where the graph model's average goes. Through
step 42 the model tracks well above prevalence; from step 43 onward it
collapses to the floor. This is the well-known dark-market shutdown in the
Elliptic timeline --- the test window contains a regime the training window
never saw.

\begin{figure}[t]
\centering
\includegraphics[width=\columnwidth]{fig3_timestep_crater.png}
\caption{Per-time-step AUC-PR on the Elliptic++ test window, 5-seed mean with
one standard deviation. Performance is strong through step 42 and collapses to
the prevalence floor from step 43, the point of a known market shutdown. A
single aggregate number hides this entirely.}
\label{fig:crater}
\end{figure}

On the procurement anchor (Table~\ref{tab:t2}) the ordering is the same and
the graph model's failure is starker: R-GCN scores $0.2808 \pm 0.0087$
against a prevalence of $0.358$ --- seed-stable, and \emph{below} the rate you
would get by guessing. The best firm-level configuration is a tree with the
publisher's own precomputed screens added ($0.4558$), which is the honest
headline for this dataset: the screens tradition contributes more than the
graph model does.

\begin{table}[t]
\caption{Mendeley EU firm-level results (within the case-control sample,
prevalence $0.358$). Lift-style reading applies. ``+screens'' rows add the
publisher's own precomputed screen columns to that baseline's inputs.}
\label{tab:t2}
\centering
\footnotesize
\setlength{\tabcolsep}{4pt}
\begin{tabular}{lll}
\toprule
Model & AUC-PR & P@18 \\
\midrule
B1 rules & 0.3426 (det.) & 0.47 \\
B2 XGBoost & 0.3925 (det.) & 0.22 \\
B3 XGBoost + graph & 0.3775 (det.) & 0.56 \\
B4 screens & 0.3811 (det.) & 0.78 \\
\textbf{B2 + screens} & \textbf{0.4558} (det.) & 0.61 \\
B3 + screens & 0.3450 (det.) & 0.39 \\
B4 + screens & 0.3874 (det.) & 0.50 \\
R-GCN focal (5 seeds) & $0.2808 \pm 0.0087$ & $0.0746 \pm 0.0363$ \\
\bottomrule
\end{tabular}
\end{table}

\subsection{Queue granularity: the unit you score is not the unit you review}
\label{sec:granularity}

Scoring transactions and scoring actors are different problems, and the
difference is measurable. On the Elliptic++ actor graph the R-GCN reaches
$0.2473$ AUC-PR against a $0.053$ prevalence --- a $4.7\times$ lift --- with an
exceptionally reliable head: precision@100 is $0.98$ and seed-stable. But
rolling those node scores up into review groups collapses the queue: 2,514
actor alerts yield precision@50 of $0.08$.

The transaction-level queue behaves in the opposite direction: 318 communities
reduce to 254 alerts with precision@50 of $0.32$, precision@100 of $0.23$ and
precision@200 of $0.135$ --- a usable triage list whose illicit coverage at the
top 50 is 7.8\%. The procurement queue gives 223 alerts from 1,283 communities
with precision@4 of $0.50$.

The practical reading is that node-level ranking quality and queue-level
precision are not interchangeable, and a paper that reports only one of them
can be describing an unusable product. We report both.

\subsection{RQ2: injection recovery separates flow motifs from award motifs}
\label{sec:injection}

Because confirmed cartels are scarce, we plant known motifs into the
unlabelled 163,327-node OCDS Georgia window and measure recovery at a fixed
budget. Table~\ref{tab:t6} is the result, and it splits cleanly.

\begin{table}[t]
\caption{Injected-member recall@2000 on the OCDS Georgia window (population
163,327; 5 seeds; unlabelled fusion mode).}
\label{tab:t6}
\centering
\small
\setlength{\tabcolsep}{3.5pt}
\begin{tabular}{lllll}
\toprule
Motif family & DOMINANT & Ensemble & Floor & GAE \\
\midrule
Coordinated cluster & $0.366$ & $\mathbf{0.923}$ & $0.124$ & $0.914$ \\
Common control & $\mathbf{0.579}$ & $0.443$ & $0.429$ & $0.466$ \\
Partition & $0.172$ & $0.034$ & $0.000$ & $0.048$ \\
Rotation & $0.084$ & $0.003$ & $0.000$ & $0.039$ \\
Cover bid & $0.007$ & $0.003$ & $0.000$ & $0.010$ \\
\bottomrule
\end{tabular}
\end{table}

Motifs that change the \emph{shape} of the graph are recovered well:
coordinated clusters at $0.923$ recall and common control at $0.579$. Motifs
that change only \emph{who wins} --- rotation, cover bidding, market partition
--- are close to invisible, at $0.003$ to $0.172$. This is a structural
limitation rather than a tuning failure: bid rotation leaves the co-bidding
topology essentially unchanged and perturbs only the award labels along it.
Detecting it needs award-sequence features, not better message passing. We
consider this one of the more useful negatives in the paper, because the
cartel literature tends to treat rotation as a flagship detectable pattern.

\subsection{RQ4: transfer is real, market-dependent, and asymmetric}

Fig.~\ref{fig:transfer} and Table~\ref{tab:t4} give the two procurement
transfer matrices. Held-out prevalences vary from $0.25$ to $0.89$, so lift
over fold prevalence is the only fair comparator.

\begin{figure}[t]
\centering
\includegraphics[width=\columnwidth]{fig4_transfer_lift.png}
\caption{Transfer lift by held-out fold. Leave-one-market-out on Garc\'ia is
positive on every market; leave-one-country-out on Mendeley averages $1.17$
but fails on its largest market. Anything at or below $1.0$ is no better than
the fold's own base rate.}
\label{fig:transfer}
\end{figure}

\begin{table}[t]
\caption{Procurement transfer matrices, 5 seeds per fold. Lift is AUC-PR
divided by the held-out fold's own prevalence; $1.0$ means no better than the
fold's base rate. Mendeley LOCO macro lift $1.17$; Garc\'ia LOMO macro lift
$1.57$.}
\label{tab:t4}
\centering
\small
\begin{tabular}{lrrr}
\toprule
Held-out fold & $n$ & AUC-PR & Lift \\
\midrule
\multicolumn{4}{l}{\textit{Mendeley EU --- leave-one-country-out}} \\
\quad country\_1 & 30 & $0.8664 \pm 0.0224$ & 1.13 \\
\quad country\_2 & 750 & $0.6464 \pm 0.0182$ & \textbf{0.90} \\
\quad country\_3 & 63 & $0.4983 \pm 0.0603$ & 1.74 \\
\quad country\_4 & 76 & $0.7427 \pm 0.0812$ & 1.03 \\
\quad country\_5 & 60 & $0.7630 \pm 0.0473$ & 1.14 \\
\quad country\_6 & 9 & $0.9527 \pm 0.0346$ & 1.07 \\
\quad country\_7 & 13 & $0.8091 \pm 0.0060$ & 1.17 \\
\midrule
\multicolumn{4}{l}{\textit{Garc\'ia Rodr\'iguez --- leave-one-market-out}} \\
\quad America & 120 & $0.3894 \pm 0.1080$ & 1.56 \\
\quad Brazil & 272 & $0.5063 \pm 0.1161$ & 1.77 \\
\quad Italy & 821 & $0.7858 \pm 0.0585$ & 1.43 \\
\quad Japan & 1665 & $0.4295 \pm 0.0223$ & 1.53 \\
\bottomrule
\end{tabular}
\end{table}

Garc\'ia leave-one-market-out is uniformly positive: every held-out market
beats its own base rate, macro lift $1.57$, and the Italy fold is strong
enough that a three-alert queue over it returns three confirmed cartels.
Mendeley leave-one-country-out is not: the macro lift of $1.17$ is carried by
small folds, and \texttt{country\_2} --- with 750 firms, by far the largest
held-out market --- comes in at $0.90$, worse than its own base rate. A macro
average over heterogeneous folds is a misleading summary here, and reporting
only that average would have converted a real limitation into a success story.

\textbf{Cross-domain.} A frozen encoder trained on one domain, probed on the
other, is asymmetric. Procurement~$\rightarrow$~financial reaches $0.1502$
against a $0.065$ prevalence, a $2.31\times$ lift.
Financial~$\rightarrow$~procurement reaches $0.2841$ against a $0.358$
prevalence --- below the base rate, i.e.\ no usable signal.

\textbf{And it is label-rich-only.} Fig.~\ref{fig:labeleff} sweeps the number
of target-domain labels available. In both directions almost every budget
below roughly a thousand labels shows a \emph{negative} transfer gain: the
frozen source probe is worse than simply training on the raw target features
you already have. Transfer here is a convenience at scale, not a cold-start
solution --- which is precisely the regime in which a practitioner would most
want it.

\begin{figure}[t]
\centering
\includegraphics[width=\columnwidth]{fig6_label_efficiency.png}
\caption{Cross-domain label efficiency. Bars are the gain of a frozen
source-domain probe over a raw-feature probe at the same label budget.
Negative almost everywhere below $k \approx 1000$: transfer does not solve the
cold-start problem in either direction.}
\label{fig:labeleff}
\end{figure}

\subsection{Robustness, ablations, and a sharp negative}
\label{sec:robustness}

Table~\ref{tab:t7} collects the component ablations. Two conventions matter.
A ``$-$'' row removes a component of the adopted model, so its delta is an
ablation cost; a ``$+$'' row adds a variant that was evaluated and
\emph{not} adopted, so its delta is a rejection margin. The two are not
comparable, and only ``$-$'' rows support a claim about the strongest
component. Every delta is measured against the reference of its own basis:
seed-0 deltas are never differenced against multi-seed means.

\begin{table}[t]
\caption{Component ablations. ``$-$'' rows are ablation costs; ``$+$'' rows
are rejection margins for variants that were measured and not adopted.}
\label{tab:t7}
\centering
\footnotesize
\setlength{\tabcolsep}{3.5pt}
\begin{tabular}{llll}
\toprule
Arm & Basis & AUC-PR & $\Delta$ \\
\midrule
GATv2-focal (headline) & 5 seeds & $0.4729 \pm 0.0525$ & ref. \\
\quad $-$ focal loss & 5 seeds & $0.4388 \pm 0.0505$ & $-0.0341$ \\
Calibrated ensemble & 5 seeds & $0.4434 \pm 0.0501$ & ref. \\
\quad $-$ unsup.\ members & 5 seeds & $0.4729 \pm 0.0525$ & $+0.0295$ \\
GATv2-focal (headline) & seed 0 & $0.5492$ & ref. \\
\quad $-$ bidir.\ edges & seed 0 & $0.3549$ & $\mathbf{-0.1943}$ \\
\quad $+$ structural chan. & seed 0 & $0.3781$ & $-0.1711$ \\
\quad $+$ line-graph chan. & seed 0 & $0.4986$ & $-0.0506$ \\
\quad $+$ context fusion & seed 0 & $0.3242$ & $-0.2250$ \\
Mendeley B2 XGBoost & det. & $0.3925$ & ref. \\
\quad $+$ dataset screens & det. & $0.4558$ & $+0.0633$ \\
\bottomrule
\end{tabular}
\end{table}

Direction-awareness is the single most valuable component: removing reverse
edges costs $0.194$ AUC-PR, independently replicating the Multi-GNN finding on
a different dataset. Focal loss is worth a real but second-order $0.034$.
Adding input channels consistently \emph{hurt} --- structural, line-graph and
context-fusion arms all scored below raw features alone --- which we read as a
capacity-versus-shift interaction rather than a claim that those channels are
uninformative in general.

\textbf{Protocol sensitivity.} Sweeping the Jaccard suppression threshold from
$0.3$ to $0.7$ and the minimum-overlap fraction across three settings leaves
the published queues unchanged on both anchors; the hit rule is likewise
robust. The queue numbers are not an artifact of a tuned deduplication
setting.

\textbf{The label-noise result.} Fig.~\ref{fig:noise} is the paper's sharpest
negative. We randomly flip a fraction of training labels and measure both
validation and test AUC-PR. Test performance \emph{rises} monotonically, from
$0.4827 \pm 0.0616$ clean to $0.5978 \pm 0.0248$ at 20\% flipped, while
validation falls from $0.9418$ to $0.5692$. The two curves move in opposite
directions.

\begin{figure}[t]
\centering
\includegraphics[width=\columnwidth]{fig5_label_noise.png}
\caption{Label-noise robustness on Elliptic++ (GATv2-focal, 3 seeds per rate).
Corrupting training labels \emph{improves} test AUC-PR while validation
collapses. Model selection on validation is optimising against the wrong
curve.}
\label{fig:noise}
\end{figure}

The explanation is that validation is drawn from the pre-shift regime and test
straddles the step-43 shift. Label noise acts as a regulariser that prevents
the model from over-committing to pre-shift structure, which costs validation
and buys test. This is our fourth independent measurement that validation
selection is blind to this shift, joining the frozen-normalisation fix, the
three added-channel degradations, and the per-step crater in
Fig.~\ref{fig:crater}.

We state plainly what this is and is not. It is a \textbf{diagnostic}, not a
recommendation: nobody should corrupt their labels. It says that under
temporal shift the standard model-selection loop is unreliable, and that a
paper reporting only validation-selected test numbers on Elliptic is reporting
something closer to luck than to skill.

\textbf{AMLworld at scale.} For completeness we report a negative at 5M edges:
tree baselines score $0.0064$ and $0.0062$ against a prevalence of $0.0104$
--- below chance. The diagnosis is structural. AMLworld's ground truth is
edge-level, and a node-first-seen split leaves test accounts with a median
degree of 3 against 13 in train, so there is almost no structure left to
score. We therefore use AMLworld in its designed role --- injection with known
ground truth --- rather than as a supervised benchmark.

% =====================================================================
\section{Explanation Study}
\label{sec:rq3}
% =====================================================================

\textbf{Choosing an attribution method.} A three-arm ablation over 50 alerts
compared GNNExplainer, PGExplainer and raw attention. PGExplainer produced
sane attributions on 49 of 50 nodes against GNNExplainer's 12 of 50, at less
than half the runtime (51.3\,s vs.\ 122.9\,s), and we adopted it. Regenerating
the bundles under PGExplainer cut implausible-fidelity explanations from 38 of
50 to 1 of 50. Attention weights, despite their popularity as a free
interpretability signal, matched GNNExplainer's weaker sanity rate.

\textbf{Motif matching and red-flag mapping.} A nine-family matcher runs over
every alert's subgraph and achieves 100\% recall on injected fixtures. On the
financial queue, 15 of 50 bundles carry a proven motif with FATF citations
attached; on the actor queue, 15 of 20. On the procurement queues the honest
answer is zero: the Mendeley alert communities and the three Garc\'ia Italy
alerts are too small --- often two firms --- to contain a shape the matcher can
prove. We report that as a coverage limitation rather than filling the gap
with a weaker claim.

\textbf{Copilot gate.} The copilot's release gate is 24 golden cases graded on
grounding and on released-output violations. It passes at 24/24 with a
grounding score of $1.0$ and zero guilt-attributing releases. Two gate
failures during development were instructive and both were fixed by prompt
rules rather than by model changes: the agent once echoed a
guilt-presupposing question in its own words while denying it, and once
reformatted a metric ($0.32 \rightarrow$ ``32\%'') in a way that broke
traceability to the artifact.

\textbf{Practitioner phase.} The instrument, rater packets and scoring rubric
are built and released, but \textbf{the human study has not been run}. We
report the protocol and treat the absence of practitioner judgements as an
explicit limitation rather than substituting author judgement for it. The
pre-registered fallback is rubric-trained graduate raters with at least one
domain expert, reported as such.

% =====================================================================
\section{A Joint Cross-Crime Extension}
\label{sec:joint}
% =====================================================================

\noindent\fbox{\parbox{0.95\columnwidth}{\small \textbf{Status: design and
protocol only.} This section specifies an experiment that has \emph{not been
run}. It contains no results, and no number in it is measured. It is included
because the design follows directly from the paper's premise and because
pre-specifying it guards against post-hoc analysis later.}}

\medskip

Everything above treats the two domains as two \emph{evaluation surfaces} for
one stack: the same architecture is trained separately on each ledger, and
\RQ{4} asks whether a representation moves between them. That framing has a
blind spot. It assumes the offender populations are disjoint --- that the
launderer and the bid rigger are different people.

They need not be. A firm that wins rigged public contracts has proceeds to
place, and the same beneficial owner may sit behind both the bidding entity
and the accounts that move the money. Under that hypothesis the two ledgers
are not two datasets but two \emph{views of one graph}, joined by
shared-control edges that neither siloed detector can see (Fig.~\ref{fig:dual}).

\begin{figure}[t]
\centering
\includegraphics[width=\columnwidth]{fig7_dual_actor.png}
\caption{The joint cross-crime design. A dual actor appears in both ledgers
under common control. A joint heterogeneous model can traverse the
shared-control edge; two siloed detectors cannot observe it at all. The
comparison is specified here and has not yet been run.}
\label{fig:dual}
\end{figure}

\subsection{Hypothesis and design}

\textbf{H1.} On actors present in both ledgers, a joint model over the merged
heterogeneous graph achieves higher recall at a fixed alert budget than the
better of two siloed detectors run independently.

\textbf{H2.} The advantage is concentrated in actors whose individual-ledger
evidence is \emph{sub-threshold} in both ledgers --- the case that motivates the
design, since actors already flagged in one ledger need no help.

\textbf{Construction.} A merged graph with typed nodes (account, entity, firm,
tender) and typed edges (transfer, co-bid, award, shared-control). An R-GCN
shares parameters across relation types, so the shared-control relation is a
first-class edge type rather than a feature. The siloed baseline is the same
architecture trained per ledger with the shared-control edges deleted, which
isolates the effect of the join rather than of model capacity.

\textbf{Ground truth.} No public dataset links a confirmed cartel firm to a
confirmed laundering account --- which is exactly why the experiment is not yet
run. The tractable design is a dual-actor injection study: plant paired motifs
(a procurement rotation clique and a financial layering chain) sharing a
controlled entity, with a controlled fraction of dual actors, and measure
recall of the planted pairs under joint versus siloed models. This inherits
the injection machinery already validated in Section~\ref{sec:injection} and
its known limitation --- award-pattern motifs are poorly recovered --- which
should be expected to bound the procurement half of any dual-actor result.

\subsection{What would falsify it}

If joint and siloed models are within noise on planted dual actors, the
cross-crime hypothesis adds nothing operationally and the siloed framing is
vindicated. Given that our measured cross-domain transfer is asymmetric,
weak in one direction and negative in the other, we regard that outcome as
genuinely plausible and would report it. The value of pre-specifying the
design is that this outcome stays publishable.

% =====================================================================
\section{Regulatory Alignment and Ethics}
% =====================================================================

The EU AI Act's Annex III places fraud detection in financial services in its
high-risk class, with obligations enforceable from 2 August 2026: technical
documentation, logging, human oversight, and information sufficient for a
deployer to interpret output. In the UAE, CBUAE Federal Decree-Law No.~6 of
2025, Article~149, requires licensed institutions to implement robust fraud
detection and to cooperate on pattern data. On the procurement side, the OECD's
2025 guidelines note that member countries spend roughly 13\% of GDP on public
procurement and that eliminating bid rigging could reduce prices by 20\% or
more.

We map obligations to shipped properties rather than to intentions. Every
alert carries a subgraph, a matched motif and a red-flag citation, which is
the interpretability obligation in concrete form. The queue is truncated to a
fixed budget, so the system is structurally a triage instrument for a human
reviewer rather than an autonomous decision-maker. Artifacts are immutable
files and the serving layer is read-only, so an audit can be run against
exactly what an analyst saw. The caveat is attached at the API boundary and
cannot be suppressed by a client, and the copilot's guilt-language guard is
enforced by a release gate rather than by prompt convention.

\textbf{Scope boundary.} The system determines no legal guilt. It ranks cases
for human investigation. Target maturity is a validated prototype (TRL 3--4)
on public and synthetic data; no personal or classified data was used at any
point. We state this in the paper for the same reason the system states it on
every response: a screening score that gets read as an accusation is a harm,
and the language is part of the design.

% =====================================================================
\section{Limitations}
% =====================================================================

\begin{itemize}
\item \textbf{The Elliptic shift dominates the financial results}, and
validation selection is blind to it --- measured four independent ways. Test
numbers on this benchmark should be read as regime-dependent.
\item \textbf{Anonymised features limit narratives.} Elliptic's features are
unnamed, so financial explanations are structural; amount-level narratives are
only available on AMLworld, which is synthetic.
\item \textbf{Mendeley is case-control.} Lift, not raw AUC-PR, is the
comparator, and population-level precision cannot be inferred from it.
\item \textbf{Garc\'ia identity coverage is four of six markets}, so
co-bidding identity graphs do not span the full benchmark.
\item \textbf{Award-pattern motifs are poorly recovered} (rotation, cover
bidding, partition), which bounds procurement claims specifically.
\item \textbf{The practitioner study's human phase has not been run}, so \RQ{3}
rests on fidelity and coverage measurements, not on practitioner judgement.
\item \textbf{The joint cross-crime experiment is unrun} (Section~\ref{sec:joint}).
\item \textbf{Compute envelope.} All published runs are CPU, full-batch.
Neighbour-sampled minibatch training on the 5M-edge AMLworld graph, and the
PNA/GIN comparisons that would need it, are outside this envelope.
\end{itemize}

% =====================================================================
\section{Reproducibility}
% =====================================================================

One YAML equals one experiment: no number in this paper exists without a
committed config. The release carries a number-to-config map with a drift
guard --- a repository test fails continuous integration if the map and the
config directory ever disagree in either direction --- plus pinned dependencies
and a Python version, per-dataset checksum manifests with licences, five
datasheets, a model card, and the tables in this paper as a regeneration
target. Cross-machine byte-reproducibility is verified for the tree baselines,
the leave-one-country-out fold and the frozen-encoder probes; torch GNN
training is deterministic same-machine only, and the $\pm 0.02$ cross-machine
variance it induces is documented and dominated by the $\pm 0.05$ seed
variance. Leakage guarantees are enforced by a dedicated test suite run as its
own continuous-integration step; the full suite is 391 tests.

% =====================================================================
\section{Conclusion}
% =====================================================================

One stack does serve both ledgers, and building it that way made the negatives
legible. Under an honest protocol, gradient-boosted trees on graph features
still beat the graph neural networks on the financial anchor, and calibration
matters more than architecture. Transfer across procurement markets is real
but market-dependent, and it fails on the largest market in the harder
benchmark --- a result that a macro average would have hidden. Cross-domain
transfer is asymmetric and only pays when labels are already plentiful.
Corrupting training labels improves test performance while validation
collapses, which says the standard selection loop is unreliable under temporal
shift.

None of those findings would have survived a single-seed, validation-selected,
aggregate-only evaluation --- which is the sense in which the protocol is the
contribution. What we ship alongside it is a screening instrument that ranks
groups for a human at a fixed budget, explains each flag in the reviewer's own
red-flag vocabulary, and refuses, by construction, to say anyone is guilty.

\section*{Acknowledgements}

[Placeholder --- funding, compute, and dataset-provider acknowledgements to be
completed before submission.]

\appendices

\section{Number-to-Config Map}
\label{app:repro}

Table~\ref{tab:repro} maps every headline number in this paper to the single
committed configuration that produces it. A repository hygiene test fails
continuous integration if this map and the configuration directory disagree in
either direction, so the mapping cannot silently rot. Each row runs as
\texttt{collusiongraph <verb> -c configs/experiment/<config>}; outputs land in
a gitignored results directory and are never committed.

\begin{table}[h]
\caption{Published number $\rightarrow$ config. Multi-seed rows run seeds 0--4
and are resumable; completed seeds are never retrained.}
\label{tab:repro}
\centering
\footnotesize
\setlength{\tabcolsep}{3pt}
\begin{tabular}{p{0.44\columnwidth}p{0.48\columnwidth}}
\toprule
Config (\texttt{configs/experiment/}) & Produces \\
\midrule
\texttt{baselines\_elliptic\_pp} & B1 $0.0576$ / B2 $0.8076$ / B3 $0.8104$ \\
\texttt{gnn\_elliptic\_pp\_gatv2\_} \texttt{focal\_multiseed} & the headline GNN number, $0.4729 \pm 0.0525$ \\
\texttt{gnn\_elliptic\_pp\_gatv2\_} \texttt{wce\_multiseed} & $-$focal ablation, $0.4388 \pm 0.0505$ \\
\texttt{gnn\_elliptic\_pp\_gatv2\_} \texttt{focal\_unidir} & $-$bidirectional edges, $0.3549$ \\
\texttt{ensemble\_elliptic\_pp\_} \texttt{multiseed} & calibrated $0.4434 \pm 0.0501$ vs.\ rank $0.0511 \pm 0.0019$ \\
\texttt{alert\_queue\_elliptic\_pp\_} \texttt{ensemble} & the financial queue: 254 alerts, P@50 $0.32$ \\
\texttt{label\_noise\_elliptic\_pp} & the noise curve, $0.4827 \rightarrow 0.5978$ \\
\texttt{sensitivity\_elliptic\_pp} & NMS / hit-rule invariance grid \\
\texttt{explainer\_ablation\_} \texttt{elliptic\_pp} & PGExplainer 49/50 vs.\ GNNExplainer 12/50 \\
\texttt{baselines\_mendeley\_eu} & B1 $0.3426$ / B2 $0.3925$ / B3 $0.3775$ / B4 $0.3811$ \\
\texttt{baselines\_mendeley\_} \texttt{screens\_ablation} & B2+screens $0.4558$ (best firm-level) \\
\texttt{gnn\_mendeley\_rgcn\_} \texttt{multiseed} & the honest negative, $0.2808 \pm 0.0087$ \\
\texttt{transfer\_loco\_matrix\_} \texttt{mendeley} & 7-fold LOCO, macro lift $1.17$ \\
\texttt{transfer\_lomo\_matrix\_} \texttt{garcia} & 4-market LOMO, macro lift $1.57$ \\
\texttt{cross\_domain\_probe\_} \texttt{proc2fin} & $0.1502$ against $0.065$ prevalence \\
\texttt{cross\_domain\_probe\_} \texttt{fin2proc} & $0.2841$ against $0.358$ prevalence \\
\texttt{label\_efficiency\_*} & both label-efficiency curves \\
\texttt{injection\_recovery\_ocds\_} \texttt{georgia\_multiseed} & coordinated cluster $0.923 \pm 0.162$ @2000 \\
\texttt{baselines\_amlworld} & the 5M-edge negative, $0.0064$ / $0.0062$ \\
\texttt{practitioner\_study} & rater packets and manifest \\
\bottomrule
\end{tabular}
\end{table}

\section{Protocol Guarantees}

Strict-inductive temporal splits, as-of feature computation,
leave-one-country/market-out entity isolation, and the alert-unit rules are
each enforced by a dedicated leakage test rather than by convention. That
suite runs standalone and as its own continuous-integration step on every
push, so a change that reintroduces leakage fails the build rather than
quietly improving a number. Cross-machine byte-reproducibility is verified for
the tree baselines, the leave-one-country-out fold and the frozen-encoder
probes. Torch GNN training is deterministic same-machine only; the resulting
$\pm 0.02$ cross-machine variance is documented and is dominated by the
$\pm 0.05$ seed variance, which is why every learned number in this paper is
reported as a multi-seed mean with its standard deviation.

\balance

\begin{thebibliography}{00}
\bibitem{ma2023} X. Ma \emph{et al.}, ``A comprehensive survey on graph anomaly detection with deep learning,'' \emph{IEEE Trans. Knowl. Data Eng.}, vol.~35, no.~12, pp.~12012--12038, 2023.
\bibitem{gadbench} J. Tang \emph{et al.}, ``GADBench: Revisiting and benchmarking supervised graph anomaly detection,'' in \emph{NeurIPS Datasets and Benchmarks}, 2023.
\bibitem{motie2024} S. Motie and B. Raahemi, ``Financial fraud detection using graph neural networks: A systematic review,'' \emph{Expert Syst. Appl.}, 2024.
\bibitem{lawal2025} A. Lawal, C. Okolie, and J. Obunadike, ``An explainable graph neural network framework for AML in cryptocurrency transactions using the Elliptic dataset,'' \emph{Int. J. Comput. Netw. Commun.}, Dec. 2025.
\bibitem{linemv2025} K. Poon \emph{et al.}, ``LineMVGNN: Anti-money laundering with line-graph-assisted multi-view graph neural networks,'' \emph{AI}, vol.~6, no.~4, art.~69, 2025.
\bibitem{multignn} E. Altman \emph{et al.}, ``Realistic synthetic financial transactions for anti-money laundering models,'' in \emph{NeurIPS}, 2023, arXiv:2306.16424.
\bibitem{elliptic} Y. Elmougy and L. Liu, ``Elliptic++: A new financial dataset for graph anomaly detection,'' arXiv:2306.06108, 2023.
\bibitem{weber2019} M. Weber \emph{et al.}, ``Anti-money laundering in Bitcoin: Experimenting with graph convolutional networks for financial forensics,'' in \emph{KDD Workshop on Anomaly Detection in Finance}, 2019.
\bibitem{garcia2022} M. J. Garc\'ia Rodr\'iguez \emph{et al.}, ``Collusion detection in public procurement auctions with machine learning algorithms,'' \emph{Autom. Constr.}, vol.~133, art.~104047, 2022.
\bibitem{mendeley2025} ``Public procurement cartels: A large-sample testing of screens using machine learning,'' \emph{Int. J. Ind. Organ.}, 2025. Dataset: Mendeley Data f3y4nrn3s6/2.
\bibitem{imhof2025} D. Imhof, O. Viklund, and M. Huber, ``Catching bid-rigging cartels with graph attention neural networks,'' arXiv:2507.12369, 2025.
\bibitem{gomes2024} L. Gomes \emph{et al.}, ``Collusion detection with graph neural networks,'' arXiv:2410.07091, 2024.
\bibitem{epjds2025} ``Detection of fraud in public procurement using data-driven methods: A systematic mapping study,'' \emph{EPJ Data Sci.}, 2025.
\bibitem{pgexplainer} D. Luo \emph{et al.}, ``Parameterized explainer for graph neural network,'' in \emph{NeurIPS}, 2020.
\bibitem{gnnexplainer} R. Ying \emph{et al.}, ``GNNExplainer: Generating explanations for graph neural networks,'' in \emph{NeurIPS}, 2019.
\bibitem{leiden} V. A. Traag, L. Waltman, and N. J. van Eck, ``From Louvain to Leiden: Guaranteeing well-connected communities,'' \emph{Sci. Rep.}, vol.~9, art.~5233, 2019.
\bibitem{gatv2} S. Brody, U. Alon, and E. Yahav, ``How attentive are graph attention networks?'' in \emph{ICLR}, 2022.
\bibitem{sage} W. Hamilton, Z. Ying, and J. Leskovec, ``Inductive representation learning on large graphs,'' in \emph{NeurIPS}, 2017.
\bibitem{rgcn} M. Schlichtkrull \emph{et al.}, ``Modeling relational data with graph convolutional networks,'' in \emph{ESWC}, 2018.
\bibitem{focal} T.-Y. Lin \emph{et al.}, ``Focal loss for dense object detection,'' in \emph{ICCV}, 2017.
\bibitem{saito2015} T. Saito and M. Rehmsmeier, ``The precision-recall plot is more informative than the ROC plot when evaluating binary classifiers on imbalanced datasets,'' \emph{PLoS ONE}, vol.~10, no.~3, 2015.
\bibitem{dominant} K. Ding \emph{et al.}, ``Deep anomaly detection on attributed networks,'' in \emph{SDM}, 2019.
\bibitem{oecd2025} OECD, \emph{Guidelines for Fighting Bid Rigging in Public Procurement (2025 Update)}, OECD Publishing, Paris, 2025, DOI 10.1787/cbe05a56-en.
\bibitem{aiact} European Union, ``Regulation (EU) 2024/1689 laying down harmonised rules on artificial intelligence (AI Act),'' \emph{Official Journal of the European Union}, 2024.
\bibitem{cbuae2025} Central Bank of the UAE, ``Federal Decree-Law No.~6 of 2025,'' Article~149, effective 16 Sep. 2025.
\bibitem{ocds} Open Contracting Partnership, ``Open Contracting Data Standard,'' \url{https://standard.open-contracting.org}.
\end{thebibliography}

\end{document}
